\subsection{导数与积分}\label{sec:distribution-derivatives-integrals}

\subsubsection{基本定义}\label{subsec:distribution-derivative-basic-definitions}

\PDEIndexMark{idx:distribution-derivatives-introduction}本节讨论分布的微分及各种应用. 讨论限于 $\mathcal{D}^{\prime}(\Omega)$ 中的分布; 完全类似的考察也适用于 $\mathcal{S}^{\prime}(\mathbb{R}^m)$. 如下定义分布的导数.

\begin{Definition}
\label{def:distribution-partial-derivative}
设 $f\in\mathcal{D}^{\prime}(\Omega)$. 那么 $f$ 关于 $x_j$ 的\textbf{导数}定义为
\begin{equation}\label{eq:distribution-partial-derivative-definition}
\left(\dfrac{\partial f}{\partial x_j},\phi\right)=-\left(f,\dfrac{\partial\phi}{\partial x_j}\right).
\end{equation}

\end{Definition}

\begin{Remark}
\label{rem:distribution-classical-derivative-agreement}
若 $f$ 属于 $C^1(\Omega)$, 则这一定义与经典导数一致, 通过分部积分即可看出. 容易看出, $\partial f/\partial x_j$ 仍属于 $\mathcal{D}^{\prime}(\Omega)$.
\end{Remark}

注意, 微分在 $\mathcal{D}^{\prime}(\Omega)$ 中是连续运算, 即读者可以证明以下结果.

\begin{Theorem}
\label{thm:distribution-derivative-continuity}
若在 $\mathcal{D}^{\prime}(\Omega)$ 中 $f_n\to f$, 则 $\partial f_n/\partial x_j\to\partial f/\partial x_j$.
\end{Theorem}

因此, 对分布而言, 交换导数与极限没有问题, 这与经典微积分中的情况形成鲜明对比.

\begin{Remark}
\label{rem:distribution-higher-derivatives}
通过递归定义高阶导数; 由定义以及试验函数的混合偏导数相等, 可明显看出混合偏导数相等. 一般地, 有
\begin{equation}\label{eq:distribution-higher-derivative-pairing}
(D^\alpha f,\phi)=(-1)^{|\alpha|}(f,D^\alpha\phi).
\end{equation}

\end{Remark}

\begin{Remark}
\label{rem:distribution-difference-quotients}
经典导数定义为差商的极限. 在某种意义下, 分布导数仍是差商的极限. 上一节通过下式定义分布的平移
\begin{equation}\label{eq:distribution-coordinate-translation-pairing}
(f(\bm{x}+h\bm{e}_j),\phi(\bm{x}))=(f(\bm{x}),\phi(\bm{x}-h\bm{e}_j)).
\end{equation}
这不一定有意义, 因为 $\bm{x}-h\bm{e}_j$ 不一定属于 $\Omega$. 然而, 对固定的 $\phi\in\mathcal{D}(\Omega)$, 只要 $h$ 充分小, $\phi(\bm{x}-h\bm{e}_j)$ 就属于 $\mathcal{D}(\Omega)$. 因而当 $h$ 很小时, \eqref{eq:distribution-coordinate-translation-pairing} 有意义, 不过 $h$ 需要多小依赖于 $\phi$. 现在得到\[
\begin{aligned}
&\lim_{h\to0}\dfrac1h\bigl[(f(\bm{x}+h\bm{e}_j),\phi(\bm{x}))-(f(\bm{x}),\phi(\bm{x}))\bigr]\\
&\quad=\lim_{h\to0}\left(f(\bm{x}),\dfrac1h(\phi(\bm{x}-h\bm{e}_j)-\phi(\bm{x}))\right)\\
&\quad=\left(f(\bm{x}),\lim_{h\to0}\dfrac1h(\phi(\bm{x}-h\bm{e}_j)-\phi(\bm{x}))\right)\\
&\quad=-\left(f,\dfrac{\partial\phi}{\partial x_j}\right)\\
&\quad=\left(\dfrac{\partial f}{\partial x_j},\phi\right).
\end{aligned}
\]
\end{Remark}
